\chapter{Các hệ thống liên quan}
\label{Chapter2}

\section{Một số hệ thống và công nghệ tương tự}

Trong những năm gần đây, cùng với sự phát triển của công nghệ thông tin trong lĩnh vực y tế, các hệ thống theo dõi sức khỏe từ xa ngày càng được quan tâm và ứng dụng rộng rãi hơn. Sự phổ biến của điện thoại thông minh, thiết bị đeo thông minh, thiết bị đo sức khỏe tại nhà và các nền tảng điện toán đám mây đã tạo điều kiện thuận lợi cho việc thu thập, lưu trữ và phân tích dữ liệu sức khỏe của người bệnh ngoài môi trường bệnh viện. Đối với các bệnh mãn tính như tăng huyết áp và đái tháo đường, việc theo dõi thường xuyên các chỉ số như huyết áp, nhịp tim, đường huyết, cân nặng hoặc các chỉ số sinh tồn khác có ý nghĩa quan trọng trong quá trình kiểm soát bệnh và phát hiện sớm dấu hiệu bất thường.

Các hệ thống Remote Patient Monitoring hiện nay có thể được triển khai theo nhiều hướng khác nhau. Một số hệ thống tập trung vào việc cung cấp thiết bị đo kết nối trực tiếp với ứng dụng di động để tự động đồng bộ dữ liệu. Một số hệ thống khác phát triển dưới dạng nền tảng quản lý bệnh mãn tính, kết hợp giữa dữ liệu đo tại nhà, huấn luyện viên sức khỏe, bác sĩ và các chương trình can thiệp hành vi. Ngoài ra, các cổng thông tin bệnh nhân cũng cho phép người bệnh xem hồ sơ sức khỏe, trao đổi với cơ sở y tế và trong một số trường hợp có thể gửi dữ liệu sức khỏe tự ghi nhận.

Trong phạm vi đề tài, nhóm tiến hành khảo sát một số hệ thống và công nghệ có liên quan đến bài toán theo dõi sức khỏe từ xa, đặc biệt là các hệ thống có liên quan đến tăng huyết áp, đái tháo đường, dữ liệu đo tại nhà, cảnh báo bất thường và hỗ trợ tương tác giữa bệnh nhân với nhân viên y tế. Các hệ thống được khảo sát bao gồm Teladoc Health/Livongo, Withings Health Solutions, Dexcom Continuous Glucose Monitoring và MyChart.

\subsection{Teladoc Health/Livongo}

\subsubsection{Giới thiệu}
\par Teladoc Health là một nền tảng chăm sóc sức khỏe từ xa cung cấp nhiều dịch vụ y tế trực tuyến, trong đó có các chương trình hỗ trợ quản lý bệnh mãn tính. Sau khi sáp nhập với Livongo, Teladoc Health mở rộng thêm các dịch vụ liên quan đến quản lý đái tháo đường, tăng huyết áp, tiền đái tháo đường và kiểm soát cân nặng. Các chương trình này thường kết hợp dữ liệu từ thiết bị kết nối, phân tích dữ liệu cá nhân hóa và sự hỗ trợ từ đội ngũ chuyên môn nhằm giúp người dùng theo dõi, hiểu và cải thiện tình trạng sức khỏe của mình. \cite{teladocChronicCare}

\par Đối với bài toán của đề tài, Teladoc Health/Livongo là một hệ thống có nhiều điểm tương đồng vì cùng hướng đến việc quản lý bệnh mãn tính từ xa. Hệ thống không chỉ ghi nhận dữ liệu sức khỏe mà còn khai thác dữ liệu đó để đưa ra phản hồi, hướng dẫn và hỗ trợ người bệnh trong quá trình kiểm soát bệnh. Đây là một mô hình tham khảo quan trọng khi xây dựng hệ thống theo dõi bệnh nhân tăng huyết áp và đái tháo đường.

\subsubsection{Những điểm nổi bật}
\begin{itemize}
    \item Hỗ trợ quản lý nhiều nhóm bệnh mãn tính phổ biến như đái tháo đường, tăng huyết áp, tiền đái tháo đường và kiểm soát cân nặng. \cite{teladocChronicCare}
    \item Có khả năng kết hợp giữa thiết bị đo sức khỏe, dữ liệu theo dõi hằng ngày và các chương trình tư vấn hoặc huấn luyện sức khỏe.
    \item Cung cấp trải nghiệm theo dõi liên tục hơn so với mô hình khám chữa bệnh truyền thống, giúp người bệnh không chỉ nhận hỗ trợ khi đến cơ sở y tế.
    \item Dữ liệu thu thập được có thể được sử dụng để cá nhân hóa khuyến nghị, hỗ trợ người dùng thay đổi hành vi và duy trì thói quen có lợi cho sức khỏe.
    \item Mô hình phù hợp với các tổ chức lớn như doanh nghiệp, công ty bảo hiểm hoặc hệ thống y tế có nhu cầu quản lý số lượng lớn người bệnh.
\end{itemize}

\subsubsection{Những điểm còn hạn chế}
\begin{itemize}
    \item Hệ thống được triển khai chủ yếu theo mô hình thương mại tại các thị trường phát triển, do đó có thể chưa phù hợp trực tiếp với điều kiện vận hành của các cơ sở y tế nhỏ hoặc các phòng khám tại Việt Nam.
    \item Một số chức năng phụ thuộc vào thiết bị đo kết nối và hạ tầng dịch vụ đi kèm, dẫn đến chi phí triển khai ban đầu có thể cao.
    \item Mô hình hoạt động thường gắn với hệ sinh thái khép kín của nhà cung cấp, gây khó khăn nếu muốn tùy biến sâu theo quy trình nghiệp vụ riêng của từng đơn vị y tế.
    \item Hệ thống hướng đến quy mô vận hành lớn, trong khi đề tài của nhóm tập trung vào mô hình thử nghiệm có phạm vi nhỏ hơn, ưu tiên khả năng triển khai đơn giản và dễ mở rộng.
\end{itemize}

\subsection{Withings Health Solutions}

\subsubsection{Giới thiệu}
\par Withings là một công ty phát triển các thiết bị sức khỏe thông minh như cân thông minh, đồng hồ theo dõi sức khỏe, nhiệt kế và máy đo huyết áp kết nối. Các thiết bị này có thể đồng bộ dữ liệu với ứng dụng di động, cho phép người dùng theo dõi các chỉ số sức khỏe theo thời gian. Bên cạnh sản phẩm cho người dùng cá nhân, Withings cũng phát triển các giải pháp dành cho tổ chức y tế, nghiên cứu lâm sàng và theo dõi bệnh nhân từ xa. \cite{withingsHealthSolutions}

\par Trong bối cảnh đề tài, Withings là một ví dụ tiêu biểu cho hướng tiếp cận sử dụng thiết bị đo sức khỏe tại nhà để thu thập dữ liệu tự động. Đặc biệt, các thiết bị đo huyết áp kết nối có thể ghi nhận huyết áp tâm thu, huyết áp tâm trương và nhịp tim, sau đó đồng bộ dữ liệu vào ứng dụng để người dùng hoặc nhân viên y tế theo dõi. Đây là hướng phát triển có thể được xem xét trong tương lai nếu hệ thống của nhóm mở rộng sang tích hợp thiết bị IoT.

\subsubsection{Những điểm nổi bật}
\begin{itemize}
    \item Cung cấp nhiều thiết bị đo sức khỏe tại nhà, trong đó có máy đo huyết áp kết nối, cân thông minh, đồng hồ thông minh và nhiệt kế thông minh. \cite{withingsHealthSolutions}
    \item Dữ liệu từ thiết bị có thể được đồng bộ với ứng dụng, giúp người dùng quan sát lịch sử đo và xu hướng thay đổi của các chỉ số sức khỏe.
    \item Giảm thao tác nhập liệu thủ công, từ đó hạn chế một phần sai sót do người dùng nhập nhầm chỉ số hoặc quên ghi nhận dữ liệu.
    \item Phù hợp với các mô hình theo dõi bệnh nhân từ xa có nhu cầu thu thập dữ liệu huyết áp, cân nặng, nhịp tim hoặc các chỉ số sinh tồn khác.
    \item Có khả năng phục vụ cả người dùng cá nhân lẫn các tổ chức cần triển khai theo dõi sức khỏe từ xa ở quy mô lớn hơn.
\end{itemize}

\subsubsection{Những điểm còn hạn chế}
\begin{itemize}
    \item Người dùng cần sở hữu thiết bị phần cứng tương thích, làm tăng chi phí sử dụng so với mô hình nhập liệu thủ công.
    \item Việc phụ thuộc vào thiết bị đo và kết nối không dây có thể phát sinh lỗi đồng bộ, lỗi pin, lỗi kết nối Bluetooth hoặc Wi-Fi.
    \item Một số thiết bị có thể chưa phổ biến tại thị trường Việt Nam, gây khó khăn cho việc triển khai đại trà trong môi trường thực tế.
    \item Hệ thống mạnh về thu thập dữ liệu từ thiết bị, nhưng khi triển khai trong một cơ sở y tế cụ thể vẫn cần bổ sung các nghiệp vụ như phân công bác sĩ, quản lý bệnh nhân, cấu hình ngưỡng cảnh báo và quy trình phản hồi cảnh báo.
\end{itemize}

\subsection{Dexcom Continuous Glucose Monitoring}

\subsubsection{Giới thiệu}
\par Dexcom là một trong những hệ thống theo dõi đường huyết liên tục phổ biến dành cho người bệnh đái tháo đường. Thay vì yêu cầu người bệnh đo đường huyết từng lần bằng cách lấy máu đầu ngón tay, hệ thống Continuous Glucose Monitoring sử dụng cảm biến đặt trên cơ thể để theo dõi đường huyết theo thời gian gần thực. Các thiết bị hiện đại có thể gửi dữ liệu đến điện thoại thông minh, hiển thị biểu đồ xu hướng và cảnh báo khi đường huyết vượt khỏi ngưỡng an toàn. \cite{dexcomCGM}

\par Đối với đề tài của nhóm, Dexcom là một hệ thống tham khảo quan trọng ở khía cạnh theo dõi đường huyết và cảnh báo bất thường. Mặc dù hệ thống của nhóm hiện chỉ đặt mục tiêu cho phép bệnh nhân nhập thủ công chỉ số đường huyết, mô hình của Dexcom cho thấy giá trị của việc theo dõi liên tục, trực quan hóa dữ liệu và cảnh báo sớm trong quản lý bệnh đái tháo đường.

\subsubsection{Những điểm nổi bật}
\begin{itemize}
    \item Cho phép theo dõi đường huyết liên tục, giúp người bệnh quan sát biến động đường huyết trong ngày thay vì chỉ dựa vào các lần đo rời rạc. \cite{dexcomCGM}
    \item Có khả năng hiển thị dữ liệu trên điện thoại thông minh, giúp người dùng dễ dàng theo dõi tình trạng hiện tại và xu hướng thay đổi.
    \item Hỗ trợ cảnh báo khi đường huyết quá cao hoặc quá thấp, góp phần giúp người bệnh phản ứng kịp thời trước các tình huống nguy hiểm.
    \item Một số giải pháp CGM cho phép chia sẻ dữ liệu với người thân hoặc nhân viên y tế, tăng khả năng hỗ trợ từ xa.
    \item Dữ liệu đường huyết liên tục có giá trị lớn trong việc đánh giá hiệu quả điều trị, chế độ ăn uống, vận động và sử dụng thuốc.
\end{itemize}

\subsubsection{Những điểm còn hạn chế}
\begin{itemize}
    \item Chi phí cảm biến và thiết bị theo dõi liên tục thường cao hơn nhiều so với mô hình nhập liệu thủ công hoặc đo đường huyết truyền thống.
    \item Hệ thống chủ yếu tập trung vào đường huyết, chưa giải quyết đầy đủ bài toán quản lý đồng thời tăng huyết áp, nhịp tim và các chỉ số sinh tồn khác.
    \item Việc sử dụng cảm biến yêu cầu người bệnh có khả năng thao tác, bảo quản và thay cảm biến đúng cách.
    \item Dữ liệu từ cảm biến cần được diễn giải cẩn thận vì quyết định điều trị vẫn phải dựa trên hướng dẫn của nhân viên y tế.
    \item Các công nghệ đo đường huyết không xâm lấn chưa được kiểm chứng đầy đủ có thể gây rủi ro nếu người bệnh sử dụng để đưa ra quyết định điều trị quan trọng. \cite{fdaGlucoseWatchWarning}
\end{itemize}

\subsection{MyChart}

\subsubsection{Giới thiệu}
\par MyChart là một cổng thông tin bệnh nhân được xây dựng trên nền tảng Epic, cho phép người bệnh truy cập một số thông tin sức khỏe cá nhân và tương tác với cơ sở y tế thông qua ứng dụng hoặc trình duyệt web. Các cổng thông tin bệnh nhân như MyChart thường hỗ trợ người dùng xem kết quả xét nghiệm, lịch hẹn, thông tin thuốc, gửi tin nhắn cho cơ sở y tế và quản lý một số thông tin liên quan đến quá trình chăm sóc sức khỏe. \cite{mychartFeatures}

\par So với các hệ thống chuyên biệt về thiết bị đo sức khỏe, MyChart đại diện cho hướng tiếp cận lấy hồ sơ sức khỏe điện tử và giao tiếp giữa bệnh nhân với cơ sở y tế làm trung tâm. Đây là một mô hình có liên quan đến đề tài vì hệ thống của nhóm cũng cần cung cấp kênh để bệnh nhân theo dõi dữ liệu cá nhân và trao đổi với bác sĩ. Tuy nhiên, đề tài tập trung cụ thể hơn vào bài toán theo dõi từ xa cho tăng huyết áp và đái tháo đường thay vì xây dựng một cổng thông tin bệnh nhân tổng quát.

\subsubsection{Những điểm nổi bật}
\begin{itemize}
    \item Hỗ trợ bệnh nhân truy cập thông tin sức khỏe cá nhân thông qua ứng dụng di động hoặc trình duyệt web. \cite{mychartFeatures}
    \item Tạo kênh giao tiếp giữa bệnh nhân và cơ sở y tế, giúp giảm phụ thuộc vào trao đổi trực tiếp tại bệnh viện.
    \item Có khả năng tích hợp với hệ thống hồ sơ sức khỏe điện tử, giúp dữ liệu bệnh nhân được quản lý tập trung hơn.
    \item Phù hợp với các cơ sở y tế có nhu cầu cung cấp dịch vụ số hóa cho người bệnh, từ đặt lịch, xem kết quả đến trao đổi thông tin chăm sóc.
    \item Mô hình cổng thông tin bệnh nhân giúp tăng tính chủ động của người bệnh trong việc theo dõi và quản lý sức khỏe cá nhân.
\end{itemize}

\subsubsection{Những điểm còn hạn chế}
\begin{itemize}
    \item MyChart là một hệ thống lớn, gắn với nền tảng hồ sơ sức khỏe điện tử của Epic, nên việc triển khai đòi hỏi hạ tầng và chi phí đáng kể.
    \item Hệ thống có phạm vi rộng, không tập trung riêng vào quy trình theo dõi hằng ngày cho bệnh nhân tăng huyết áp và đái tháo đường.
    \item Việc cấu hình chức năng theo dõi từ xa, thu thập dữ liệu tại nhà hoặc cảnh báo phụ thuộc vào cách từng cơ sở y tế triển khai.
    \item Đối với các nhóm phát triển nhỏ hoặc các cơ sở y tế muốn thử nghiệm mô hình RPM ở quy mô ban đầu, một hệ thống tổng quát như MyChart có thể phức tạp hơn nhu cầu thực tế.
\end{itemize}

\section{Những vấn đề còn tồn đọng}

\par Nhìn chung, các hệ thống và công nghệ kể trên đều cho thấy tiềm năng lớn của việc ứng dụng công nghệ thông tin vào theo dõi và quản lý sức khỏe từ xa. Tuy nhiên, khi đặt trong phạm vi bài toán xây dựng một hệ thống Remote Patient Monitoring cho người bệnh tăng huyết áp và đái tháo đường tại Việt Nam, vẫn còn một số vấn đề đáng chú ý.

\par Thứ nhất, nhiều hệ thống hiện nay phụ thuộc khá nhiều vào thiết bị phần cứng chuyên dụng. Các thiết bị như máy đo huyết áp kết nối, cân thông minh, cảm biến đường huyết liên tục hoặc đồng hồ thông minh giúp giảm thao tác nhập liệu và tăng độ liên tục của dữ liệu, nhưng đồng thời làm tăng chi phí triển khai. Với nhiều bệnh nhân, đặc biệt là người lớn tuổi hoặc người có điều kiện kinh tế hạn chế, việc yêu cầu sở hữu thiết bị chuyên dụng có thể trở thành rào cản lớn. Do đó, mô hình nhập liệu thủ công vẫn có ý nghĩa trong giai đoạn đầu vì dễ tiếp cận hơn, không phụ thuộc vào thiết bị và phù hợp với mục tiêu xây dựng mô hình thử nghiệm của đề tài.

\par Thứ hai, các hệ thống thương mại lớn thường được thiết kế cho thị trường, quy trình và hạ tầng y tế của các quốc gia phát triển. Khi áp dụng vào môi trường khác, đặc biệt là tại các cơ sở y tế nhỏ hoặc phòng khám, hệ thống có thể cần được điều chỉnh nhiều để phù hợp với quy trình quản lý bệnh nhân, phân công bác sĩ, vai trò của y tá, cách tiếp nhận cảnh báo và phương thức trao đổi với người bệnh. Đây là lý do đề tài của nhóm tập trung xây dựng các vai trò như bệnh nhân, bác sĩ, y tá và quản trị viên để phản ánh rõ hơn nhu cầu vận hành trong một hệ thống quản lý bệnh nhân.

\par Thứ ba, mỗi hệ thống thường chỉ mạnh ở một khía cạnh nhất định. Các nền tảng như Dexcom rất mạnh trong theo dõi đường huyết liên tục nhưng không giải quyết toàn diện bài toán tăng huyết áp. Các thiết bị của Withings hỗ trợ tốt việc đo và đồng bộ dữ liệu sức khỏe tại nhà nhưng vẫn cần một lớp phần mềm nghiệp vụ để quản lý phân công, cảnh báo và xử lý dữ liệu theo từng bệnh nhân. Các cổng thông tin bệnh nhân như MyChart mạnh về truy cập hồ sơ và giao tiếp với cơ sở y tế nhưng không nhất thiết tối ưu cho bài toán cảnh báo theo ngưỡng hằng ngày của một nhóm bệnh cụ thể.

\par Thứ tư, dữ liệu sức khỏe là loại dữ liệu nhạy cảm, đòi hỏi hệ thống phải chú trọng đến bảo mật, phân quyền và kiểm soát truy cập. Trong các hệ thống theo dõi từ xa, dữ liệu có thể đến từ nhiều nguồn như bệnh nhân nhập thủ công, thiết bị đo tại nhà, ứng dụng di động hoặc nhân viên y tế. Nếu không có cơ chế phân quyền phù hợp, bác sĩ hoặc y tá có thể truy cập nhầm dữ liệu của bệnh nhân không thuộc phạm vi phụ trách, hoặc bệnh nhân có thể xem dữ liệu không phải của mình. Vì vậy, hệ thống cần thiết kế rõ ràng các quyền truy cập theo vai trò và theo quan hệ phân công.

\par Thứ năm, cơ chế cảnh báo trong hệ thống RPM cần được thiết kế cẩn thận. Nếu ngưỡng cảnh báo quá nhạy, hệ thống có thể tạo ra quá nhiều cảnh báo không cần thiết, khiến bác sĩ và y tá khó ưu tiên xử lý. Ngược lại, nếu ngưỡng cảnh báo quá rộng, hệ thống có thể bỏ sót các dấu hiệu bất thường quan trọng. Trong phạm vi đề tài, nhóm lựa chọn cơ chế cảnh báo dựa trên ngưỡng cố định hoặc ngưỡng do bác sĩ cấu hình, chưa áp dụng học máy hoặc mô hình dự đoán rủi ro. Cách tiếp cận này phù hợp với giai đoạn thử nghiệm vì đơn giản, dễ giải thích và dễ kiểm soát.

\par Cuối cùng, việc theo dõi sức khỏe từ xa không thể thay thế hoàn toàn quá trình khám chữa bệnh trực tiếp. Dữ liệu đo tại nhà có thể chịu ảnh hưởng bởi cách đo, thời điểm đo, thiết bị đo, thói quen nhập liệu và tình trạng thực tế của bệnh nhân. Do đó, hệ thống được xây dựng với mục tiêu hỗ trợ theo dõi, phát hiện sớm dấu hiệu bất thường và cung cấp thêm dữ liệu tham khảo cho nhân viên y tế, chứ không đưa ra chẩn đoán hoặc quyết định điều trị thay cho bác sĩ. Đây cũng là định hướng quan trọng để đảm bảo tính khả thi và an toàn của đề tài trong phạm vi đồ án tốt nghiệp.